\chapter{Logging, tracing, and observability}
\label{ch:observability}

\begin{aside}
An interactive program that does not log is opaque when things
go wrong. \maya{} ships a small observability toolkit: a ring
buffer for last-N messages, structured tracing for the render
pipeline, and a debug overlay you can toggle at runtime. None
of it is mandatory; all of it pays off the first time a user
files a bug.
\end{aside}

\section{Why not just std::cerr?}

A TUI app cannot use \code{std::cerr} freely because stderr is
the user's screen. Writing log messages to stderr corrupts the
display.

Three workable alternatives:

\begin{itemize}[leftmargin=*]
  \item Redirect stderr to a file at startup.
  \item Buffer log messages in memory; surface on demand.
  \item Write to a separate log file always.
\end{itemize}

\maya{} supports all three. The default is the ring buffer:
in-memory by default, dump on request.

\section{The Logger interface}

\file{maya/include/maya/app/context.hpp} exposes:

\begin{lstlisting}
class Logger {
public:
    void debug(std::string_view msg);
    void info(std::string_view msg);
    void warn(std::string_view msg);
    void error(std::string_view msg);

    template <class... Args>
    void info(std::string_view fmt, Args&&... args) {
        info(std::format(fmt, std::forward<Args>(args)...));
    }

    std::vector<LogEntry> recent(size_t n) const;
    void dump_to(std::ostream& os) const;
};
\end{lstlisting}

Four levels (debug, info, warn, error). Each takes a message;
the \code{std::format}-aware overloads support typed
arguments. \code{recent(n)} returns the last $n$ entries;
\code{dump\_to} writes them to a stream.

\subsection*{LogEntry}

\begin{lstlisting}
struct LogEntry {
    std::chrono::system_clock::time_point time;
    Level level;
    std::string category;
    std::string message;
};
\end{lstlisting}

Time, level, category (subsystem tag), message. Categories let
you filter logs by source.

\section{The ring buffer}

The default logger is a fixed-size ring:

\begin{lstlisting}
class RingLogger {
    std::vector<LogEntry> buffer_;
    size_t head_ = 0;
    size_t count_ = 0;
    mutable std::mutex mu_;

public:
    explicit RingLogger(size_t capacity);
    void append(LogEntry e);
    std::vector<LogEntry> snapshot() const;
};
\end{lstlisting}

Capacity is typically 1024 or 4096 entries. Older entries are
overwritten in place; the ring stays at constant memory.

The mutex is fine-grained: a single mutex protects the buffer.
Logging from multiple threads is supported (worker tasks, the
main thread).

\subsection*{Snapshot}

\begin{lstlisting}
std::vector<LogEntry> RingLogger::snapshot() const {
    std::lock_guard lk{mu_};
    std::vector<LogEntry> out;
    out.reserve(count_);
    for (size_t i = 0; i < count_; ++i) {
        out.push_back(buffer_[(head_ - count_ + i) % buffer_.size()]);
    }
    return out;
}
\end{lstlisting}

Returns the entries in chronological order. Threadsafe; can be
called while logging continues.

\section{Structured logging}

For tracing, \maya{} uses a slightly richer entry shape:

\begin{lstlisting}
struct TraceEntry {
    std::chrono::steady_clock::time_point time;
    std::string event;
    std::unordered_map<std::string, std::string> attributes;
};
\end{lstlisting}

The attributes map lets you attach context: \code{frame=42},
\code{node\_count=128}, \code{cell\_count=1920}. When you
later inspect the trace, you can filter and group by attribute.

\subsection*{Tracing the render pipeline}

The pipeline emits trace events at each stage:

\begin{lstlisting}
trace(\"view.begin\", {{\"frame\", std::to_string(frame_no)}});
auto e = P::view(model);
trace(\"view.end\", {{\"frame\", std::to_string(frame_no)},
                    {\"node_count\", std::to_string(count_nodes(e))}});
\end{lstlisting}

A trace pair: \code{begin}/\code{end} with the same key
attributes. The duration is the difference; the attribute set
lets you correlate.

\section{The debug overlay}

A widget that renders the most recent log entries:

\begin{lstlisting}
Element render_log_overlay(const Logger& l) {
    auto entries = l.recent(20);
    std::vector<Element> lines;
    for (auto& e : entries) {
        auto style = level_style(e.level);
        lines.push_back(h(
            text(format_time(e.time)) | Dim,
            text(\" \"),
            text(e.message) | style
        ));
    }
    return box(v(lines)) | border<Single> | pad<1>;
}
\end{lstlisting}

Toggled with a hotkey (typically F12 or Ctrl-Shift-L). The
overlay sits on top of the app via the z-stack from
Chapter~\ref{ch:elements}.

\section{Categories and filters}

Tag log messages with a category:

\begin{lstlisting}
logger.info(\"render\", \"frame {} ({} cells dirty)\", n, dirty);
logger.info(\"network\", \"GET {} -> 200\", url);
logger.info(\"input\", \"key {} pressed\", key.name());
\end{lstlisting}

Then filter:

\begin{lstlisting}
auto render_logs = logger.recent_filtered(
    100, [](auto& e) { return e.category == \"render\"; });
\end{lstlisting}

For development, set an environment variable to control which
categories are recorded:

\begin{lstlisting}
MAYA_LOG_CATEGORIES=render,network ./my_app
\end{lstlisting}

Only the listed categories are appended; others are dropped
cheaply.

\section{Telemetry to a remote sink}

For deployed apps, log to a file or service. The pattern:

\begin{lstlisting}
auto telemetry_cmd = Cmd<Msg>::task([entries] {
    upload(entries);
    return TelemetryDone{};
});
\end{lstlisting}

A background task uploads batches. The main loop is unaffected;
the user does not see the upload.

\subsection*{Privacy}

Be careful about what you log. PII (personally identifiable
information), file paths, secrets --- redact before recording.
A simple redactor:

\begin{lstlisting}
std::string redact(std::string_view msg) {
    // Replace anything that looks like an email, a token, a path.
    return std::regex_replace(std::string{msg},
        std::regex{R\"(\\w+@\\w+\\.\\w+)\"}, \"[REDACTED]\");
}
\end{lstlisting}

The list of patterns is app-specific. Err on the side of
redacting too much.

\section{Crash diagnostics}

If your app crashes, what survives?

\subsection*{The ring buffer is volatile}

A SIGSEGV does not flush. You need a signal handler:

\begin{lstlisting}
extern \"C\" void crash_handler(int sig) {
    logger.dump_to(std::ofstream{\"/tmp/myapp.crashlog\"});
    std::signal(sig, SIG_DFL);
    std::raise(sig);
}
\end{lstlisting}

Install with \code{std::signal(SIGSEGV, crash\_handler)}. The
handler must be async-signal-safe; in practice ``open a file
and write'' is OK but allocating memory is not. \maya{}'s
\code{dump\_to} is designed to be async-signal-safe.

\subsection*{Symbolicated stacks}

The crash log should include a stack trace. Use
\code{std::stacktrace} (C++23):

\begin{lstlisting}
auto trace = std::stacktrace::current();
for (auto& frame : trace) {
    log_line(std::format(\"{}\", frame));
}
\end{lstlisting}

Compile with debug symbols (\code{-g}); the trace will name
functions and line numbers.

\section{Performance considerations}

Logging is on the hot path. Cheap logging:

\begin{itemize}[leftmargin=*]
  \item String formatting is the slowest part; do it lazily.
  \item Disabled levels (e.g. \code{debug} in release) should
    cost zero --- one branch on a constant flag, not even a
    function call.
  \item Lock contention is real; if multiple threads log
    heavily, the buffer mutex becomes a bottleneck. Use
    thread-local buffers and a periodic merge.
\end{itemize}

The simplest lazy pattern:

\begin{lstlisting}
if (logger.level_enabled(Level::Debug)) {
    logger.debug(\"value = {}\", expensive_compute());
}
\end{lstlisting}

The expensive compute only runs if debug is on.

\section{Pitfalls}

\begin{warning}
\textbf{Logging in the inner loop.} Even cheap logging adds up
at 10\,000 calls per frame. Sample (log every 100th occurrence)
or gate by level.
\end{warning}

\begin{warning}
\textbf{Logging secrets.} A bug report with a config file
containing API keys is a security incident. Redact aggressively.
\end{warning}

\begin{warning}
\textbf{Lost messages on crash.} A buffer that lives only in
memory disappears with the process. Pair with a crash handler
that flushes.
\end{warning}

\begin{warning}
\textbf{Overlapping logs.} Multiple threads writing to the
same buffer can produce interleaved messages. The mutex
serialises them, but if you log a multi-line message, the
lines may not be adjacent. Log one message per call.
\end{warning}

\section*{Workshop}

\begin{enumerate}[leftmargin=*]
  \item Add a ring logger to a small app. Log on every key
    press. Toggle the debug overlay; verify entries appear.
  \item Install a crash handler that dumps the log on
    \code{SIGSEGV}. Force a crash (deref a null pointer);
    verify the log is written.
  \item Add structured tracing to the render pipeline (begin
    + end events with attributes). Compute frame durations
    from the trace.
\end{enumerate}

\begin{recap}
A TUI cannot use stderr freely; use a ring buffer or a log
file. Structured tracing captures context with attributes.
A debug overlay surfaces recent entries at runtime. A crash
handler dumps the buffer on SIGSEGV. Categories and levels
let you filter cheaply. Redact secrets before they leave the
process.
\end{recap}
